\section{Frontier Workshop Notes}

\subsection{Workshop focus areas}

El docente también recomendó organizar workshops periódicos, para concentrarse en resolver los open problems dentro de offline RL:

\begin{itemize}
    \item Lab session: Replay buffer introspection, revisión manual de high-risk trajectories.
    \item Governance: simulate counterfactual audits + compliance docs in same meeting.
    \item Tooling: swap data versioning backends, evaluate time-to-recover metrics.
\end{itemize}

\begin{figure}[H]
\centering
\includegraphics[page=9,width=0.92\textwidth]{07_cs224r_offline_rl_2025.pdf}
\caption{La diapositiva 9 resume la agenda del workshop y el checkpoint\protect\footnotemark}
\end{figure}
\footnotetext{La diapositiva 9 enumera el checkpoint del workshop y a los responsables.}

\subsection{Resumen de la sección}

El modelo de workshop puede traducir rápidamente los research insights en operational improvements, evitando los knowledge silo.

